\fancyhead[L]{\small BERICHTIGUNG DER SCHNEIDEN EINER WAAGE.}
\SourceLeaf{517}\pdfbookmark[0]{Berichtigung der Schneiden einer Waage}{schneiden_waage}
\begin{center}
{\Large [BERICHTIGUNG DER SCHNEIDEN EINER WAAGE.]\par}
\vspace{7mm}\hrule\vspace{2mm}
Göttingische gelehrte Anzeigen. 1837, 13. März.
\vspace{2mm}\hrule\vspace{5mm}
\end{center}
In der Sitzung der Königlichen Gesellschaft der Wissenschaften vom 28. Januar nahm Hofrat \textsc{Gauss} den Vortrag von Professor \textsc{Weber}, über den im 22. Stück dieser Zeitschrift berichtet wurde, zum Anlass, über einen eng verwandten Gegenstand vorzutragen. Wir geben hier den wesentlichen Inhalt seines Vortrags wieder.

Es geht um ein neues Justierverfahren zur Erfüllung einer wesent\-lichen Bedingung bei feineren Hebelwaagen, deren Bedeutung bisher nicht hinreichend gewürdigt worden zu sein scheint. Solche Waagen haben drei prismatische Schneiden. Die mittlere, nach unten gerichtete Schneide ruht auf einem harten horizontalen Lager aus Stein oder Stahl und dient bei der Bewegung des Waagebalkens als Drehachse. Die beiden anderen an den Enden des Waagebalkens sind nach oben gerichtet. Auf jeder von ihnen schwebt das Tragestück, an dem die Waagschale hängt. Die Tragestücke selbst bestehen aus gehärtetem Stahl. Ihre unteren, auf den Schneiden aufliegenden Flächen sind vollkommen eben und hochglanzpoliert.

Eine wesentliche Bedingung ist nun, dass die beiden äußeren Schneiden zur mittleren parallel sein sollen. In der Tat. da man die Waage vor jedem Austausch der Gewichte in einer Schale zunächst arretiert und dabei das Tragestück von der Schneide abhebt, kann man nie damit rechnen, dass das Tragestück nach dem Lösen der Arretierung wieder \textit{genau} so auf der Schneide aufliegt wie zuvor. Das ist unschädlich, wenn die betreffende Schneide zur mittleren parallel ist, führt aber zu einem veränderten Moment, wenn die Schneiden voneinander abweichen. Eine unvollkommene Justierung in dieser Hinsicht ist eine Hauptursache dafür, dass bei oft wiederholten Wägungen die Ergebnisse zuweilen erheblich stärker voneinander abweichen, als man
\SourceLeaf{518}es angesichts der ausgezeichneten Verarbeitung und Empfindlichkeit einer Waage sonst erwarten sollte.

Mit den Mitteln, die Instrumentenbauer bisher gewöhnlich zur parallelen Ausrichtung der Schneiden eingesetzt haben, lässt sich nicht die ganze wünschenswerte Genauigkeit erzielen. Bei feinen Waagen wie bei astronomischen Instrumenten obliegt zudem die feinste Justierung demjenigen, der die Waage benutzt. Vom Hersteller ist sie nicht zu verlangen.

Das Verfahren, das Hofrat \textsc{Gauss} mit bestem Erfolg für diese Justierung eingesetzt hat, beruht auf den folgenden Überlegungen.

Bei den Schwingungen des Waagebalkens ändert die zu prüfende äußere Schneide ihre Lage im Raum. Ist sie zur (ruhenden) mittleren Schneide parallel, so sind alle diese Lagen zueinander parallel. Anders verhält es sich, wenn die äußere Schneide zur mittleren nicht parallel ist. Nehmen wir zur Veranschaulichung an, dass die äußere und die mittlere Schneide in derselben Ebene liegen, ihre Richtungen aber vom Beobachter weg auseinanderlaufen. In diesem Fall bewegt sich die äußere Schneide beim Schwingen des Waagebalkens auf einer Kegelfläche. Ihr vom Beobachter abgewandtes Ende steigt oder sinkt gegenüber dem näheren Ende, je nachdem, ob der Hebelarm, an dem sich die Schneide befindet, steigt oder sinkt. Dasselbe gilt für das Tragestück, das die Schneide ständig berührt.

Welcher der beiden Fälle vorliegt, lässt sich erkennen, wenn man auf dem Tragestück einen Planspiegel befestigt. Am vorteilhaftesten ist es, den Spiegel so anzubringen, dass seine Ebene annähernd senkrecht zur Schneide steht, wobei man es damit nicht zu genau nehmen muss. Im ersten der beiden Fälle bleibt der Spiegel während der Bewegung des Waagebalkens zu sich selbst parallel. Im zweiten Fall gilt das nicht. Das Bild eines Gegenstands, der sich in geeigneter Entfernung vor dem Spiegel befindet, bleibt also im ersten Fall unverrückt. Im zweiten Fall steigt oder sinkt es dagegen, wie man leicht einsieht, mit dem betreffenden Hebelarm. Das Umgekehrte würde eintreten, wenn die beiden Schneiden vom Beobachter weg zusammenliefen. Dann würde das Bild beim Steigen des Waagebalkenarms sinken und umgekehrt.

Bei einem sehr guten Spiegel lässt sich selbst eine
\SourceLeaf{519}äußerst kleine Verschiebung des Bildes mit einem Fernrohr sicher und scharf erkennen. Hofrat \textsc{Gauss} verwendete als Gegenstand eine in Millimeter eingeteilte Skala, die etwa 5 Meter vor dem Spiegel senkrecht aufgestellt war. Das 35-fach vergrößernde Fernrohr stand in annähernd derselben Entfernung. So erschien das Bild eines Millimeters etwa 20 Bogensekunden groß, wovon man noch Zehntel schätzen kann. Solange die Schneide noch nicht vollkommen justiert war, bewegte sich das Bild der Skala am Fadenkreuz des Fernrohrs ganz regelmäßig auf und ab, während der Waagebalken schwang.

Für mathematisch gebildete Leser genügt der Hinweis, dass sich auf diese Weise erkennen lässt, nach welcher Seite die Schneiden voneinander abweichen und, mit hinreichender Genauigkeit, wie groß die Abweichung ist. Kennt man außerdem die Steigung der Gewinde der Justierschrauben, kann man die Justierung sicher durchführen.

Der Vollständigkeit halber seien hier noch ein paar weitere Umstände erwähnt.

Verwendet man einen etwas größeren Spiegel (der von Hofrat \textsc{Gauss} verwendete Spiegel war 75 Millimeter hoch und mit einer besonderen Vorrichtung auf dem Tragestück befestigt), müssen die Schalen mit hinreichend schweren Gewichten belastet werden. Sonst würde das Tragestück seitlich umkippen.

Oben wurde vorausgesetzt, dass die zu prüfende äußere Schneide mit der mittleren in derselben Ebene liegt. Hat man die mittlere genau horizontal eingestellt und steht der Waagebalken horizontal, so liegt auch die äußere Schneide horizontal und kann allenfalls seitlich abweichen. In der Regel trifft diese Voraussetzung jedoch nicht vollkommen genau zu. Die äußere Schneide kann in dieser Stellung etwas geneigt sein, oder eines ihrer Enden kann etwas höher liegen als das andere. Bei dem beschriebenen Prüfverfahren erkennt man dies daran, dass sich das Spiegelbild beim Steigen des Waagebalkenarms zugleich seitwärts und beim Sinken zur entgegengesetzten Seite bewegt. Dabei ist jedoch zu beachten, dass dieser Fehler, wenn er bei einer Waage eines geschickten Instrumentenbauers vorhanden ist, jedenfalls viel zu klein sein wird, um einen merklichen Fehler in den Wägeergebnissen hervorzurufen. Daher gibt es selbst bei den besten Waagen keine Vorrichtungen, mit denen sich dieser Anteil der Abweichung von der Parallelität beseitigen lässt.

